\chapter{Self Hosted Setup}

\section{Einleitung (Quellen fehlen)}

In den vergangenen Jahren ist in der digitalen Welt eine zunehmende Abhängigkeit von Cloud-Diensten und abonnementbasierten Online-Dienstleistungen zu beobachten \cite{statista_cloud_2025}. 

Gleichzeitig werden digitale Dienste und Daten nicht nur für wirtschaftliche Zwecke (wie Werbung) genutzt, sondern sind auch sicherheitsrelevant. Private Unternehmen und staatliche Institutionen verfügen über umfangreiche Möglichkeiten zur Erhebung, Verarbeitung und Analyse von Nutzerdaten, mit denen das Verhalten von Nutzern gezielt beeinflusst werden kann. Diese Einflussnahme bezieht sich nicht nur auf personalisierte Werbung, sondern betrifft zunehmend auch politische Meinungsbildung und die Inhalte, die den Nutzern auf sozialen Medien ausgespielt werden. 

Je mehr Daten über eine Person gesammelt werden, desto präziser lassen sich individuelle Profile erstellen, die einem digitalen Abbild, einem sogenannten "digitalen Zwilling", nahekommen. Diese Entwicklung wirft erhebliche datenschutzrechtliche und ethische Fragen auf. 

In den USA schafft beispielsweise der Cloud Act die rechtliche Grundlage dafür, dass Behörden unter bestimmten Voraussetzungen auf Daten von US-Unternehmen zugreifen können, auch wenn diese außerhalb der Vereinigten Staaten gespeichert sind. 


In der Europäischen Union werden Maßnahmen diskutiert und teilweise umgesetzt, die den Zugriff auf digitale Kommunikationsinhalte betreffen. Im Rahmen von Initiativen zur Bekämpfung von illegalen Inhalten, wie etwa der Verbreitung von Missbrauchsdarstellungen, wird unter anderem über verpflichtende Prüfmechanismen für Kommunikationsdienste ("Chatkontrolle") debattiert. Diese Ansätze stehen im Spannungsfeld zwischen Sicherheitsinteressen und dem Schutz der Privatsphäre, da sie potenziell auch verschlüsselte Kommunikation betreffen könnten.

Auch aus technischer Sicht bestehen Einschränkungen: Zwar setzen viele Dienste auf Ende-zu-Ende-Verschlüsselung, jedoch verbleiben durch Metadaten, serverseitige Komponenten und die Abhängigkeit von zentralen Infrastrukturen weiterhin potenzielle Angriffs- und Zugriffspunkte. Zudem stellen zentralisierte Plattformen ein attraktives Ziel für Angriffe dar, was regelmäßig zu Datenlecks und Sicherheitsvorfällen führt. Insbesondere im Kontext digitaler Vermögenswerte wie Kryptowährungen können solche Vorfälle erhebliche finanzielle Schäden verursachen.

Vor diesem Hintergrund wächst bei sicherheits- und datenschutzbewussten Nutzern der Bedarf nach Lösungen, die eine höhere Kontrolle über die eigenen Daten ermöglichen, um Abhängigkeiten von externen Dienstanbietern zu reduzieren und die Verarbeitung sensibler Daten eigenständig zu kontrollieren.

Ausgehend von den oben genannten Faktoren wurde im Rahmen dieses Projekts eine sichere, vollständig selbst gehostete Infrastruktur für einen Mandanten konzipiert. Bei dem Mandanten handelt es sich um eine vermögende Privatperson mit Schwerpunkt im Bereich Kryptowährungen, für die Diskretion, Privatsphäre und Schutz vor Überwachung von zentraler Bedeutung sind. Die physische Sicherheit der eingesetzten Systeme wird für diese Arbeit als gegeben vorausgesetzt, sodass  der Fokus auf der Bereitstellung und Absicherung der Software und des Netzwerkes liegt.

Ziel dieser Arbeit ist die Entwicklung eines wartbaren und hochsicheren Systems, dass realistische Bedrohungsszenarien wie Malware, Supply-Chain-Angriffe sowie lokale Angriffen standhält und gleichzeitig eine datensouveräne und benutzerfreundliche Nutzung zentraler Dienste ermöglicht. 

